\begin{abstract}
\noindent
\textbf{Question.} Can reduced-parameter monitoring provide reliable Water Quality Index (WQI) prediction for data-scarce river systems?

\textbf{Data.} 32,968 river monitoring visits (888 stations, 87 rivers, 2018--2021) from the CPCB National Water Quality Monitoring Programme were cleaned, joined and audited. Only 14.7\% of visits report all eight parameters needed for the CPCB weighted-arithmetic WQI, and all of them fall in a single year (2020). By contrast, 75.8\% report pH, dissolved oxygen (DO) and BOD.

\textbf{Methods.} The capped full-panel WQI of 4,861 fully observed visits (593 stations) is the fixed target. Eight model families, from the no-ML renormalised formula to XGBoost and Gaussian processes, are trained on twelve named parameter subsets and on all 255 possible subsets. They are evaluated on held-out stations, rivers, states and seasons. Further experiments cover parameter importance (eight methods), data scarcity (2--100\% of training stations), conformal and Gaussian-process prediction intervals, missing data, measurement noise, and sensitivity to the index definition. The model is then deployed on 20,132 visits that lack a full panel.

\textbf{Results.} An exact variance decomposition shows that the CPCB unit weights ($\propto 1/S_i$) concentrate 87.5\% of index variance in pH and BOD, while fecal coliform, total coliform and EC together explain under 0.3\%. All model-based importance methods agree. At unseen stations, XGBoost on pH, DO and BOD reaches \R{}\,=\,0.957, RMSE 2.9 points and 90.6\% category accuracy, against \R{}\,=\,0.863 for re-computing the index from the same three parameters. pH, BOD and turbidity reach 0.977, and adding DO gives 0.994. No reduced-parameter prediction is ever off by more than one category. Accuracy holds for unseen rivers and states (\R{}\,$\ge$\,0.94). Learned models need data from about 25 or more stations to beat the formula, and GPR is best when data are scarcest. Split-conformal intervals keep nominal coverage (90\% intervals of about $\pm$4.7 points for the three-parameter model). Field-level pH noise ($\sigma=0.1$) costs more accuracy than dropping five parameters. Conclusions are specific to the index definition: under equal weights the same three parameters explain only 62\% of the variance. Deployment yields the first full-panel WQI estimates for 255 stations. It also reveals a BOD reporting floor outside 2020 that confounds year-to-year comparisons.

\textbf{Conclusion.} For the CPCB index, a precisely measured three-parameter panel (pH, DO, BOD), plus a turbidity probe where possible, gives reliable WQI estimates with calibrated uncertainty at unmonitored-panel sites. Fecal coliform exceedance should be reported alongside it, because the CPCB weighting ignores it.
\end{abstract}
